\documentclass[12pt]{article}
\usepackage{../../ficha}

\begin{document}

\ficha{Tabelas de Dupla Entrada, Diagramas em Árvore e Lei dos Grandes Números}

\instrucoes{Apresenta as probabilidades na forma de fração irredutível. Usa a regra de Laplace: $P(A) = \dfrac{\text{n.º de casos favoráveis}}{\text{n.º de casos possíveis}}$, quando os casos são equiprováveis.}

\secao{Tabelas de Dupla Entrada}{Lançam-se dois dados cúbicos equilibrados, com as faces numeradas de 1 a 6, e somam-se os números das faces voltadas para cima.}
\begin{exercicios}
  \item Constrói uma tabela de dupla entrada com todas as somas possíveis.
  \item Quantos resultados possíveis existem?
  \item Qual é a soma mais provável? Qual é a sua probabilidade?
  \item Qual é a probabilidade de a soma ser 7?
  \item Qual é a probabilidade de a soma ser maior do que 9?
  \item Qual é a probabilidade de a soma ser um número par?
  \item Qual é a probabilidade de sair o mesmo número nos dois dados?
  \item Qual é a probabilidade de a soma ser 1? E menor do que 13? Como se chamam estes acontecimentos?
\end{exercicios}

\secao{Produto e Diferença}{Lançam-se dois dados e calcula-se o produto dos números saídos.}
\begin{exercicios}
  \item Constrói a tabela de dupla entrada com todos os produtos.
  \item Qual é a probabilidade de o produto ser 12?
  \item Qual é a probabilidade de o produto ser ímpar?
  \item Qual é a probabilidade de o produto ser maior do que 20?
  \item Qual é a probabilidade de o produto ser múltiplo de 5?
  \item Se, em vez do produto, se calcular a diferença (maior menos menor), qual é a probabilidade de a diferença ser 0? E 2?
\end{exercicios}

\secao{Diagramas em Árvore}{Resolve usando um diagrama em árvore.}
\begin{exercicios}
  \item Lança-se uma moeda três vezes. Constrói o diagrama em árvore e indica todos os resultados possíveis.
  \item Com a moeda lançada três vezes, qual é a probabilidade de sair exatamente duas caras?
  \item Qual é a probabilidade de sair pelo menos uma cara?
  \item Um restaurante oferece 3 entradas, 4 pratos principais e 2 sobremesas. Quantas refeições diferentes (entrada, prato, sobremesa) se podem compor?
  \item A Inês tem 3 camisolas (azul, branca, verde) e 2 calças (preta, ganga). Escolhe ao acaso uma camisola e umas calças. Qual é a probabilidade de escolher a camisola azul com as calças de ganga?
  \item Um saco tem 2 bolas vermelhas e 1 azul. Tiram-se duas bolas, uma após a outra, \textbf{sem reposição}. Qual é a probabilidade de ambas serem vermelhas?
  \item Repete o exercício anterior, mas \textbf{com reposição}.
\end{exercicios}

\secao{Frequência Relativa e Lei dos Grandes Números}{Resolve.}
\begin{exercicios}
  \item Lançou-se um dado 600 vezes e saiu o número 6 em 112 lançamentos. Qual é a frequência relativa do acontecimento «sair 6»? Compara-a com a probabilidade teórica.
  \item Lançou-se uma moeda 20 vezes e saiu cara 14 vezes. Pode-se concluir que a moeda é viciada? Justifica.
  \item Uma pionese caiu com o bico para cima 324 vezes em 900 lançamentos. Estima a probabilidade de cair com o bico para cima.
  \item O que diz a lei dos grandes números?
  \item Num saco há bolas brancas e pretas. Retirou-se uma bola, com reposição, 500 vezes, e saiu bola branca 301 vezes. Sabendo que há 20 bolas no saco, estima quantas são brancas.
\end{exercicios}

\secao{Desafios}{Resolve os problemas.}
\begin{exercicios}
  \item Lançam-se dois dados. Qual é a probabilidade de pelo menos um deles mostrar 6?
  \item Uma caixa tem 3 bolas numeradas de 1 a 3. Tiram-se duas bolas, sem reposição, e forma-se um número de dois algarismos (a primeira é o algarismo das dezenas). Qual é a probabilidade de o número ser par?
  \item O Rui e a Ana jogam: lançam dois dados; se a soma for par ganha o Rui, se for ímpar ganha a Ana. O jogo é justo? Justifica.
  \item Num jogo, ganhas se, ao lançar dois dados, a diferença dos números for 1 ou 2. Qual é a probabilidade de ganhar?
  \item Numa família com três filhos, qual é a probabilidade de serem todos do mesmo sexo? (Considera os dois sexos igualmente prováveis.)
\end{exercicios}

\end{document}
